\movedhead{What the null for H10 actually is}
\label{sec:h10null}\label{app:h10null}
H10 is a conditional hypothesis, and the condition does work that a naive
count ignores. The rule selects a family on having a non-zero effect at
\texttt{low} and then asks whether a second, independent measurement of the
same quantity is smaller. A family qualifies partly because its true effect
is large and partly because its noise at \texttt{low} happened to point away
from zero, and the second measurement does not inherit that luck. So the
qualifying cases shrink on re-measurement at a rate above one half even when
effort does nothing at all, which is the standard regression artefact of
selecting on a noisy baseline \citep{barnett2005regression}. Reporting
$2^{-n}$ against a count of $n$ successes treats the rate as one half and
overstates the evidence.

The rate can be measured rather than assumed. In
\texttt{analysis/h10\us placebo.py} the effort label is shuffled among the
responses to the same item within a model, over 400 shuffles, and the
registered H10 rule including its bootstrap filter is run on each shuffled
replicate. The rule, the sample sizes and the item set are untouched; only
the labels move, so no real effect of effort can survive and whatever the
rule still finds is the artefact alone. Using the seven confirmatory
families, a qualifying case shrinks with probability 0.654, not 0.500, and
the running estimate is stable from the first hundred shuffles. Across all
nine families the rate is 0.706. The rate is further from one half in the
inclusive set because the two exploratory families have the smallest true
effects, so a larger share of what qualifies them is noise.

The matched null event is the one H10 asserts: every qualifying case shrinks.
Its measured probability is 0.0125 on the confirmatory families and 0.0200
across all nine. Against the coin flip's $2^{-9}$ for the nine confirmatory
cases observed, one in five hundred and twelve, the measured figure is larger
by a factor of six. The observed result is therefore unlikely under the null,
but by well under an order of magnitude less than a count alone suggests, and
the paper claims the measured figure.

Two features of the measurement are worth stating because they cut in
opposite directions. The shuffled replicates qualify more cases than the real
data does, 11.48 on average with a range of 6 to 20 against the observed 9,
so the joint event is being evaluated over a longer run of cases than the one
actually reported; that makes 0.0125 smaller than the figure a nine-case
comparison would give. Raising the measured per-case rate to the ninth power
gives 0.022, which is the matched figure if qualifying cases were
independent. They are not, since cases within a model share the same
responses, so the truth is above the directly measured 0.0125 and the paper
uses that bound rather than the smaller number. Either way the conclusion is
the same: the pattern is unlikely under the null by roughly one part in fifty
to one part in eighty, and not one part in five hundred.
